\chapter{Conclusion et perspectives}
\label{chap:conclusion}

\section{Bilan}

Ce projet tutoré a permis de concevoir LogiFlow, un TMS routier cohérent sur le plan
architectural : domaine explicite (commande $\rightarrow$ dossier $\rightarrow$ voyage),
monolithe modulaire hexagonal, SPA Angular filtrée par rôle, agents IA isolés derrière un
BFF Spring, et infrastructure AWS reproductible (une EC2, Terraform, Ansible, SSM) sans SSH exposé.

Les diagrammes de ce rapport synthétisent les décisions structurantes : isolation de Flask,
validation JWT Keycloak, graphe Modulith, topologie VPC / Caddy / Compose, et pipeline
GHCR $\rightarrow$ SSM.

\section{Apports pour l'équipe}

\begin{itemize}[leftmargin=*]
  \item Conception métier et glossaire partagé.
  \item Architecture logicielle (DDD léger, hexagonal, Modulith).
  \item Full-stack (Angular + Spring + PostgreSQL).
  \item Intégration d'agents IA (Flask, LLM cloud, OSRM, outils Spring) avec contraintes de sécurité.
  \item Infrastructure as code et culture DevSecOps (OIDC GitHub, SSM, scans CI).
\end{itemize}

\section{Difficultés rencontrées}

Plusieurs frictions ont structuré nos arbitrages. Brancher correctement Keycloak
(OIDC + PKCE) de bout en bout --- frontend, issuer public, JWKS interne, profils Spring
\texttt{dev}/prod --- a demandé une configuration stricte pour que le local reproduise le
cloud. Côté IA, la latence d'un LLM local a conduit à adopter un fournisseur cloud et à
séparer calcul déterministe / rédaction ; le streaming SSE a dû être préservé de bout en
bout (Caddy \texttt{flush\_interval}, Gunicorn gthread). Sur l'infrastructure, le budget
étudiant a motivé le choix d'une seule EC2 sans NAT, au prix d'une surface réseau différente
d'une architecture multi-tiers classique --- documentée dans les ADR. Enfin, la coordination
de quatre dépôts (front, back, IA, infra) a nécessité des contrats d'API stables et des
revues régulières avec l'encadrant.

\section{Perspectives}

\begin{enumerate}[leftmargin=*]
  \item Activer MFA Keycloak de bout en bout sur tous les comptes cloud.
  \item Renforcer Spring $\leftrightarrow$ Flask (mTLS ou OAuth2 client-credentials).
  \item Enrichir le copilote (outils, RAG) et la planification (contraintes métier avancées).
  \item Extraire un module (\texttt{planning} ou \texttt{tracking}) en service séparé
        \emph{seulement} si la charge le justifie --- les frontières \texttt{api} sont déjà prêtes.
  \item Industrialiser le monitoring (métriques, alertes) au-delà des healthchecks et budgets AWS.
\end{enumerate}

LogiFlow constitue ainsi une base solide pour poursuivre le développement et une
industrialisation progressive.
